
\section*{How long until the gas reaches the casing?}
For the sustained nominal load, the independent moving-wall calculation
exposes the casing after about 104--111 s in the sustainer and 138--148 s in
the booster. The clock starts when hot gas first reaches bare phenolic.
Burnthrough occurs when the moving char surface reaches the casing.
\begin{table}[H]\centering
\small\begin{tabular}{lrrrl}\toprule
Motor & Liner [mm] & Radial [s] & ANSYS [s] & Earlier radial casing event \\\midrule
Sustainer Al & 4.76 & 110.5 & 110.6 & All Al liquid at 99.1 s \\
Sustainer CF & 4.76 & 104.3 & 104.3 & Casing at 1000 $^\circ$C at 95.1 s \\
Booster Al & 6.35 & 147.6 & --- & All Al liquid at 139.4 s \\
Booster CF & 6.35 & 137.9 & --- & Casing at 1000 $^\circ$C at 129.7 s \\
\bottomrule\end{tabular}
\end{table}
The calculation applies 2504 K gas continuously at 800 psi with an initial
film coefficient of 1000 W/(m$^2$ K), gas blowing and zero hot-side radiation.
The Al comparison follows pure-metal fusion and retains the molten metal in
place. The CF comparison follows oxygen-free T700/M21 resin decomposition
and holds the residue in place. The study has no measured Iris wall heat flux
or fired-liner recession. All four times therefore describe the stated thermal
scenarios. The casing can lose pressure capacity while phenolic remains.

\paragraph{What most affects the estimate.}
Halving and doubling the nominal film coefficient gives 57--207 s for
the sustainers and 75--276 s for the boosters with the evolving casings.
Section~\ref{sec:film} gives all twelve times and the heat flux entering the
liner. With nominal heating, changing only the energy required to remove a
kilogram of char gives 72--185 s for the aluminium sustainer
(Section~\ref{sec:removal}). A measured firing heat-flux history
and liner recession would constrain the physical estimate.

\paragraph{Verification.}
The four nominal burnthrough times come from a conservative radial finite-volume
solver. A second finite-element calculation from ignition gives endpoints within
0.073 s for all four cases (Section~\ref{sec:sectorcheck}). Both use the same
gas exposure and material inputs. Native ANSYS verifies the final char layer
in both CF cases with a prescribed casing temperature (Section~\ref{sec:nativecfend}).
Section~\ref{sec:whole} gives the full native runs that passed audit. A dash marks a pending full native run.
\clearpage
